\section{Introduction}

Every multi-agent orchestration topology carries a hidden bet about the tasks it will face. A hierarchical SOP pipeline assumes that work decomposes into clean, sequential subtasks. A flat conversation loop assumes that progress emerges through iterative refinement. A DAG-based workflow assumes that subtasks can be identified upfront and executed in parallel along explicit data dependencies. These structural assumptions get baked in at design time, yet nobody states them, measures them, or tests them against actual target tasks. When the assumptions match, the topology thrives. When they don't, it fails in ways that practitioners blame on model limitations rather than architectural misfit.

The field has been asking the wrong question. ``Which multi-agent framework is best?'' treats topology as a fixed choice, like picking a database engine. The right question is ``which topology fits \emph{this} task?'' --- and answering it requires understanding what each topology assumes and what each task demands.

Consider a concrete example from our experiments. A practitioner deploying hierarchical SOP on a debugging task that requires iterative hypothesis-revision ($I{=}0.95$) would observe 14\% accuracy: the planner decomposes the problem into per-function analyses, but race conditions are cross-function interactions that resist decomposition. Switching to debate on the same task yields 77\% accuracy, because adversarial critique surfaces the cross-cutting bugs that modular analysis misses. The 63-percentage-point gap is not a model problem. It is a topology problem, and it is entirely predictable from the task's structural profile.

At least eighteen multi-agent LLM frameworks now exist (AutoGen \cite{wu_2023}, MetaGPT \cite{hong_2023}, CAMEL \cite{li_2023}, ChatDev \cite{qian_2024}, DyLAN \cite{liu_2024}, AFlow \cite{aflow_2025}, and others), each committing to a different topology and each evaluating only its own design on its own tasks. The framework literature and the benchmark literature developed in parallel with zero cross-cluster citation edges \cite{mohammadi_2025}. Practitioners choose topologies based on hearsay, single-framework ablations, or marketing claims that confound topology effects with differences in LLM backbone, tool access, and prompt engineering. The cost is concrete: wasted API tokens, inflated latency, and failed deployments that a better-matched topology would have handled.

The closest precedent comes from multi-agent debate. Smit et al.~\cite{smit_2024} systematically benchmarked debate strategies and found that most do not consistently outperform chain-of-thought baselines \cite{wei_2022}. This is widely cited as evidence against multi-agent coordination. We argue it reflects something more specific: the debate topology's structural prior (high iterativeness, adversarial refinement) was tested on tasks that did not benefit from adversarial back-and-forth. The problem was not that multi-agent coordination fails generally; the wrong topology was applied to the wrong task. What is missing is a framework that makes these structural assumptions explicit and testable.

We propose that topology effectiveness becomes predictable once task structure is measured along three dimensions: \textbf{decomposability} (D), the degree to which a task separates into independently solvable subtasks; \textbf{iterativeness} (I), the fraction of steps requiring revision of earlier outputs; and \textbf{tool diversity} (T), the number of distinct tool categories the task demands. Each topology occupies a characteristic region in this D-I-T space. The alignment hypothesis follows directly: the topology whose structural prior best matches a task's D-I-T coordinates will outperform alternatives on that task.

We test this through \textbf{OrchestraBench}, a controlled benchmarking framework that isolates the topology variable. OrchestraBench implements three canonical topologies (flat conversation, hierarchical SOP, and multi-agent debate) as lightweight wrappers sharing an identical LLM backbone (Claude Opus 4.6), identical tool access, and identical token budget. We evaluate on 82 custom tasks (12 easy, 70 hard) spanning coding, reasoning, and research categories, with a secondary backbone check on Claude Sonnet 4.6. Each task is annotated with D, I, and T values, enabling direct measurement of topology-task alignment.

This paper makes three contributions:

\begin{enumerate}
\item
  \textbf{A formal taxonomy of orchestration topologies with D-I-T task characterization.} We define three evaluated topology classes and formalize each topology's structural prior as a position in D-I-T space. Two additional topologies (role-playing, RL-orchestrated) are characterized but not yet evaluated, demonstrating the framework's extensibility. The three task-structure dimensions achieve inter-annotator agreement averaging 0.84 Cohen's $\kappa$.
\item
  \textbf{OrchestraBench: controlled cross-topology evaluation at scale.} We provide the first head-to-head comparison of three orchestration topologies on 82 custom tasks under strictly controlled conditions (same LLM, same tools, same budget), yielding 794 total runs across three random seeds, two backbones, and an agent-count control experiment. The hypothesis is falsifiable: if topology effectiveness were unrelated to task structure, alignment distance would show no correlation with performance ranking.
\item
  \textbf{The first topology selection guide grounded in task characteristics.} A lightweight DIT routing classifier selects the best topology with 90\% accuracy (95\% bootstrap CI: [84\%, 95\%]), against a 33\% random baseline. We distill these predictions into actionable selection rules mapping D-I-T profiles to recommended topologies.
\end{enumerate}
